\documentclass[11pt]{article}
\usepackage[T1]{fontenc}
\usepackage[utf8]{inputenc}
\usepackage{lmodern,amsmath,amssymb,amsthm,mathtools,microtype,xcolor}
\usepackage[margin=1in]{geometry}
\usepackage[colorlinks=true,linkcolor=blue!50!black,urlcolor=blue!50!black,citecolor=blue!50!black]{hyperref}
\newtheorem{theorem}{Theorem}[section]
\newtheorem{proposition}[theorem]{Proposition}
\newtheorem{lemma}[theorem]{Lemma}
\newtheorem{corollary}[theorem]{Corollary}
\theoremstyle{remark}\newtheorem{remark}[theorem]{Remark}
\newcommand{\E}{\mathbb E}\newcommand{\Prb}{\mathbb P}\newcommand{\R}{\mathbb R}
\hypersetup{pdftitle={Bayesian circuit learning at mean-field scaling},pdfauthor={ChatGPT-6 Astra; Kaarel Hanni; Dmitry Vaintrob}}
\setlength{\emergencystretch}{3em}
\title{Bayesian circuit learning at mean-field scaling:\\a variational principle and width-independent generalization}
\author{ChatGPT-6 Astra \and Kaarel H\"anni \and Dmitry Vaintrob\thanks{Authors are listed in alphabetical order.}}
\date{October 15, 2026}
\begin{document}
\maketitle
\begin{abstract}
We prove polynomial sample complexity for Bayesian learning of Boolean functions
by fixed-depth tanh networks at mean-field scaling. The Gaussian prior has
first-layer standard deviation $d^{-1/2}$, internal-layer standard deviation
$N^{-1/2}$, and readout standard deviation $N^{-1}$. It is first restricted to
networks with absolute margin greater than $1/2$ on the entire Boolean cube;
learning then conditions on noiseless labels. For every sufficiently large
width, a sample bound polynomial in the input dimension and the target's
bounded-weight network complexity ensures correct-label posterior probability
strictly greater than $2/3$ on almost all test inputs. The sample bound is
independent of width; the required width threshold is not quantified.
The proof identifies the exact finite-domain large-deviation rate, shows that
full-margin minimizers have quadratic feature tilts with polynomially bounded
entrywise coefficients, and realizes their feature covariances by polynomial
finite networks. We also prove a target-mass bound for every fixed depth and
all $\gamma\geq1$ in the strong-prior family, and a shallow-network learning
theorem valid at every width throughout that family.
\end{abstract}

\section{Model and conclusions}
Let $d\geq1$, $L\geq2$, $\mathcal X=\{-1,1\}^d$, and
$\bar x=(1,x_1,\ldots,x_d)$. For a bounded $1$-Lipschitz activation
$\phi:\R\to[-1,1]$, a depth-$L$, hidden-width-$N$ network is
\[
 h^0(x)=\bar x,\qquad h^s(x)=\phi(W^s h^{s-1}(x))\ (1\leq s<L),
 \qquad z_w(x)=W^Lh^{L-1}(x).
\]
There are no independent hidden biases beyond the constant input coordinate.
For a Boolean function $f$ define
\[
 A_{N,f}=\{w:f(x)z_w(x)>1/2\text{ for all }x\in\mathcal X\},
 \qquad A_N=\bigcup_f A_{N,f}.
\]
Let $C_L(f)$ be the least integer $C\geq1$ for which a width-$C$ network
implements $f$ with strict signed output margin $1/2$ on the entire
Boolean cube and all its weights bounded in absolute value by $C$.
We assume this complexity is finite for the target.

Let $\mu_{N,\gamma}$ be the product Gaussian law with standard deviations
\[
 \sigma_1=N^{(1-\gamma)/2}/\sqrt d,\quad
 \sigma_s=N^{-\gamma/2}\ (2\leq s<L),\quad
 \sigma_L=N^{-(\gamma+1)/2}.
\]
Define the weighted norm and the induced Boolean prior by
\[
 \kappa_N(w)=\frac dN\|W^1\|_F^2+
 \sum_{s=2}^{L-1}\|W^s\|_F^2+N\|W^L\|_2^2,
 \qquad
 \pi_{N,\gamma}(f)=\frac{\mu_{N,\gamma}(A_{N,f})}{\mu_{N,\gamma}(A_N)}.
\]
For $x_1,\ldots,x_n$ drawn independently from any distribution $D$ on
$\mathcal X$, the posterior is this prior conditioned on
$g(x_i)=f(x_i)$ for all $i$. Write
\[
 R_D(\Pi)=\E_{g\sim\Pi}\Prb_{x\sim D}\{g(x)\ne f(x)\}.
\]
The normalization by $\mu_{N,\gamma}(A_N)$ cancels from every prior ratio used
below. The global margin restriction is essential to this model description;
we do not replace it by training-set margins.

\begin{theorem}[Mean-field learning at sufficiently large width]\label{thm:critical-learning}
Fix depth $L\geq2$ and activation $\phi=\tanh$. There is a polynomial $p_L$
with the following property. For every target $f$ with $C=C_L(f)<\infty$
and every $\varepsilon\in(0,1)$, there is a width threshold
$N_0=N_0(d,L,f,\varepsilon)\geq C$ such that for every $N\geq N_0$,
every $\delta\in(0,1)$, every input distribution $D$, and every
\[
 n\geq\frac6\varepsilon\left[p_L(d,C)+\log(1/\delta)\right],
\]
with probability at least $1-\delta$ over the training sample $S_n$,
\[
 R_D(\Pi_{N,1})\leq\varepsilon/3,
 \qquad
 \Prb_{x\sim D}\!\left\{
 \Pi_{N,1}(g(x)=f(x)\mid S_n)>2/3\right\}\geq1-\varepsilon.
\]
Here fixed depth means that both the coefficients and degree of $p_L$
may depend on $L$. The polynomial is independent of width. No polynomial
bound on $N_0$, or
all-width version of this deep-network statement, is asserted.
\end{theorem}

\begin{theorem}[Uniform shallow-network learning]\label{thm:shallow}
For $L=2$ there is a polynomial $p$ such that the following holds for every
$\gamma\geq1$, every target $f$ with $C=C_2(f)$, every $N\geq C$, every
$\varepsilon,\delta\in(0,1)$, and every input distribution $D$.
If
\[
 n\geq\frac{2}{\varepsilon}
 \left[p\bigl(d,C,\log(2/\varepsilon)\bigr)+\log(1/\delta)\right],
\]
then, with probability at least $1-\delta$ over the training sample,
$R_D(\Pi_{N,\gamma})\leq\varepsilon$.
The same posterior-majority classifier has error at most $2\varepsilon$.
The polynomial is independent of both $N$ and $\gamma$.
\end{theorem}
\begin{corollary}[Confident correct prediction]
Under the hypotheses of Theorem~\ref{thm:shallow}, if
\[
 n\geq\frac{6}{\varepsilon}
 \left[p\bigl(d,C,\log(6/\varepsilon)\bigr)+\log(1/\delta)\right],
\]
then with training-sample probability at least $1-\delta$,
\[
 \Prb_{x\sim D}\!\left\{
   \Pi_{N,\gamma}(g(x)=f(x)\mid S_n)>\frac23\right\}
 \geq1-\varepsilon.
\]
Thus an input counts as successful only when the posterior assigns its correct
label probability strictly greater than $2/3$.
\end{corollary}
\begin{proof}
On every input failing the displayed confidence condition, the posterior
error probability is at least $1/3$. Consequently the probability of
confidence failure is at most $3R_D(\Pi)$. Apply
Theorem~\ref{thm:shallow} with risk tolerance $\varepsilon/3$.
\end{proof}
The first theorem treats the infinite-width mean-field regime through a
statement valid for all sufficiently large finite widths. The second theorem
is stronger in its width quantifier for shallow networks and requires only a
bounded Lipschitz activation. The proof first establishes a common target-mass
bound. An exact large-deviation representation then identifies the critical
rate, and a new forward construction compresses its full-margin minimizers.
A finite relevant-function argument converts the resulting count into the
learning guarantee. The separate shallow argument gives its all-width result
without a limiting variational analysis.

% Standalone section: macros deliberately limited to standard AMS commands.
\section{A cylinder construction valid throughout \texorpdfstring{$\gamma\geq1$}{gamma >= 1}}
The following construction bounds the target's prior cost at every width $N\geq C$,
without a width threshold diverging as $\gamma\downarrow1$.

\begin{proposition}[Uniform cloning mass bound]
Fix $L\geq2$ and a bounded $1$-Lipschitz activation. There is a function $Q_L(d,C)$ bounded by a polynomial such that for every $\gamma\geq1$, every Boolean target with
$C_L(f)\leq C$, and every $N\geq C$,
\[
 -\frac{2}{N^\gamma}\log\Pr\{f(x)z_w(x)>1/2\text{ for all }x\}
 \leq Q_L(d,C).
\]
The prior variances and absence of later-layer biases are exactly those in the
main model.
\end{proposition}
\begin{proof}
Choose a width-$C$ teacher, increase its readout by a factor of three, and write
its hidden matrices as $B^s$ and readout as $v$. Thus the teacher has margin
strictly greater than $3/2$, hidden entries bounded by $C$, and readout entries
bounded by $3C$. If the original witness has smaller width, first use its actual
width in the construction below; all estimates only improve.
Partition $\{1,\ldots,N\}$ into $C$ groups of sizes $k_a$ differing by at most
one. Then
\[
 1\leq k_a,\qquad \frac{\max k_a}{\min k_a}\leq2,
 \qquad \frac N{k_a}\leq2C.
\]
Let $u_a=k_a^{-1/2}{\bf1}_{\text{group }a}$ and let $P$ be the orthogonal
projection onto their span. The exact clone $w_*$ has first-layer row $B^1_b$
for a neuron in group $b$, hidden entry $B^s_{ba}/k_a$ from group $a$ to group
$b$, and readout $v_b/k_b$ on group $b$. Every hidden activation of the exact
clone is group-constant. Moreover,
\[
 \|W_*^s\|_{\mathrm{op}}\leq\sqrt2 C^2\quad(2\leq s<L),
 \qquad \sqrt N\|W_*^L\|_2\leq3\sqrt2 C^2.
\]
The first bound follows by restricting to the group-constant subspace: the
resulting matrix has entries $\sqrt{k_b/k_a}B^s_{ba}$, and its Frobenius norm
is at most $\sqrt2 C^2$.

For fixed $\gamma\geq1$, set $t=t(N)=N^{(1-\gamma)/2}\leq1$.
This is the prescribed width-dependent Gaussian prefactor; it is constant
in $N$ only at $\gamma=1$. The random coordinates
$\sqrt d W^1_{ij}$, $\sqrt N W^s_{ij}$ for $2\leq s<L$, and
$\alpha_j=NW^L_j$ are independent centered Gaussians of standard deviation
$t$. Impose the following cylinder event $\mathcal C_\rho$, for
$0<\rho\leq1$:
\begin{itemize}
\item every $\sqrt d W^1_{ij}$ is within $\rho t$ of
$\sqrt d (W_*^1)_{ij}$;
\item for every hidden row $j$ and group $a$, the coordinate
$\langle\sqrt N W^s_{j,\cdot},u_a\rangle$ is within $\rho t$ of its
cloned value $\sqrt{N/k_a}B^s_{ba}$;
\item every $\alpha_j$ is within $\rho t$ of its cloned value $Nv_b/k_b$.
\end{itemize}
All these constraints concern independent scalar Gaussian coordinates.
The orthogonal coordinates of each internal row remain independent prior
Gaussians. Consequently write
\[
 W^s=W_*^s+R^s+E^s,\qquad R^s=R^sP,\quad E^sP=0.
\]
On the cylinder, $\|R^s\|_{\mathrm{op}}\leq\|R^s\|_F\leq\rho t\sqrt C$.
The free matrix is distributed as $(t/\sqrt N)G^s(I-P)$, with $G^s$ a standard
Gaussian matrix, independently of the cylinder coordinates. Let
\[
 b_L=2+\sqrt{2\log(2L)}.
\]
The Gaussian operator-norm bound and a union bound imply, conditionally on
$\mathcal C_\rho$, that all $\|E^s\|_{\mathrm{op}}\leq b_Lt$ with probability
at least $1/2$. This is also valid when $L=2$, where there are no such matrices.

Let $e_s(x)=N^{-1/2}\|h^s_w(x)-h^s_*(x)\|_2$. The input constraints give
$e_1(x)\leq\rho(d+1)/\sqrt d\leq\rho(d+1)$. Since $E^sh^{s-1}_*=0$,
boundedness and Lipschitz continuity give, simultaneously for every input,
\[
 e_s(x)\leq (\sqrt2 C^2+\sqrt C+b_L)e_{s-1}(x)+\rho\sqrt C.
\]
Put $A=2+3C^2+b_L$. Induction yields
$e_s(x)\leq\rho(d+L+1)A^{L-2}$ for $s<L$. At the output,
\[
 |z_w(x)-z_*(x)|
 \leq3\sqrt2 C^2 e_{L-1}(x)+\rho.
\]
Hence choosing
\[
 \rho=\bigl[2+9C^2(d+L+1)A^{L-2}\bigr]^{-1}
\]
ensures output discrepancy at most $1/2$. All networks in the cylinder with
the operator-norm event therefore implement $f$ with margin greater than one.

It remains to price the cylinder. Its number of constrained coordinates is
\[
 m=N\bigl[d+2+(L-2)C\bigr].
\]
The sum of their squared centers is at most $NB_0$, where one may take
\[
 B_0=d(d+1)C^2+2(L-2)C^4+36C^4.
\]
Indeed, a first-layer center has magnitude at most $\sqrt d C$, an internal
center at most $\sqrt{2C}C$, and a readout center at most $6C^2$.
For $Z\sim N(0,t^2)$ and $|Z-a|\leq\rho t$,
\[
 \Pr\{|Z-a|\leq\rho t\}
 \geq\frac{2\rho}{\sqrt{2\pi}}
       \exp\{-a^2/t^2-\rho^2\}.
\]
Multiplication and $t^{-2}=N^{\gamma-1}$ show
\[
 -\log\Pr(\mathcal C_\rho)
 \leq N^\gamma B_0
 +N[d+2+(L-2)C]\log\frac{e\sqrt{2\pi}}{2\rho}.
\]
The desired probability is at least $\Pr(\mathcal C_\rho)/2$. Since
$N\leq N^\gamma$ and $N^\gamma\geq1$, this proves the claim with
\[
 Q_L(d,C)=2B_0+\frac{4[d+2+(L-2)C]}{\rho}+2.
\]
Here we used $\log(e\sqrt{2\pi}/(2\rho))\leq2/\rho$ for $0<\rho\leq1$.
The displayed expression is a polynomial in $d,C$ at every fixed depth.
\end{proof}


\section{An exact variational principle at the critical exponent}
\label{sec:critical-variational}

For arbitrary fixed depth, the critical prior admits a direct finite-domain
large-deviation description. This result requires neither a propagation-of-chaos
assumption nor differentiability of the activation. It provides the variational
foundation for the compression argument that follows.

Fix $\gamma=1$, enumerate $\mathcal X=\{x^1,\ldots,x^m\}$, where $m=2^d$,
and put $r=L-1$. All limits in this section hold with $d$ and $L$ fixed.
Write $E=[-1,1]^m$ and let $\mathcal P(E)$ carry its weak topology.
For $\nu\in\mathcal P(E)$ define
\[
 K(\nu)=\int_E hh^{\mathsf T}\,d\nu(h),\qquad
 Q_0(a,b)=\frac{\bar x^a\cdot\bar x^b}{d}.
\]
For any positive semidefinite $m\times m$ matrix $Q$, let
\[
 \mathsf T_Q(u)=\bigl(\phi((Q^{1/2}u)_a)\bigr)_{a=1}^m,
 \qquad
 \mathsf M_Q=(\mathsf T_Q)_\#\mathcal N(0,I_m).
\]
Here $T_\#\rho$ is the distribution of $T(U)$ when $U$ has distribution $\rho$,
and $Q^{1/2}$ is the positive semidefinite square root. Relative entropy is
\[
 D(\nu\Vert\rho)=
 \begin{cases}
 \displaystyle\int\log\frac{d\nu}{d\rho}\,d\nu,&\nu\ll\rho,\\
 +\infty,&\text{otherwise}.
 \end{cases}
\]
Finally define the Gaussian energy
\[
 \mathcal G_Q(z)=
 \begin{cases}
 \frac12 z^{\mathsf T}Q^+z,&z\in\operatorname{range}(Q),\\
 +\infty,&z\notin\operatorname{range}(Q),
 \end{cases}
\]
where $Q^+$ is the Moore--Penrose inverse.

\begin{theorem}[Critical finite-domain large deviations]
\label{thm:critical-ldp}
Let $Z_N=(z_w(x^a))_{a=1}^m$ under the unconditioned prior
$\mu_{N,1}$. The sequence $Z_N$ satisfies a large-deviation principle
with speed $N$ and good rate function
\begin{equation}\label{eq:critical-rate}
 \begin{split}
 I(z)=\inf_{\nu_1,\ldots,\nu_r\in\mathcal P(E)}
 \bigg\{&D(\nu_1\Vert\mathsf M_{Q_0})
  +\sum_{s=2}^{r}D\bigl(\nu_s\Vert\mathsf M_{K(\nu_{s-1})}\bigr)
  +\mathcal G_{K(\nu_r)}(z)\bigg\}.
 \end{split}
\end{equation}
In particular, for every open $O\subseteq\mathbb R^m$ and closed
$F\subseteq\mathbb R^m$,
\[
 -\inf_{z\in O}I(z)
 \leq\liminf_{N\to\infty}\frac1N\log\Prb(Z_N\in O),
 \qquad
 \limsup_{N\to\infty}\frac1N\log\Prb(Z_N\in F)
 \leq-\inf_{z\in F}I(z).
\]
The infimum in \eqref{eq:critical-rate} is attained whenever it is finite.
No covariance matrix in the statement is required to be nonsingular.
\end{theorem}

\begin{proof}
We give an exact representation of the finite-width law and then apply
Sanov's theorem and continuous contraction~\cite{dembozeitouni}.

Let $\Gamma_m=\mathcal N(0,I_m)$. Independently for each
$s=1,\ldots,r$, let $U_{s,1},\ldots,U_{s,N}$ have law $\Gamma_m$, and set
\[
 \eta_{s,N}=\frac1N\sum_{i=1}^N\delta_{U_{s,i}}.
\]
Also let $V_N=G/\sqrt N$, where $G\sim\Gamma_m$ is independent of these
empirical measures. Recursively define
\[
 \nu_{1,N}=(\mathsf T_{Q_0})_\#\eta_{1,N},\qquad
 \nu_{s,N}=(\mathsf T_{K(\nu_{s-1,N})})_\#\eta_{s,N}\quad(2\leq s\leq r).
\]
These are exactly the joint laws of the network's empirical hidden-activation
vectors on $\mathcal X$. Indeed, first-layer preactivation vectors are
independent Gaussians with covariance $Q_0$. Conditional on the preceding
layer, each subsequent row of weights gives an independent Gaussian vector
with covariance
\[
 \frac1N\sum_{i=1}^N
 \bigl(h_i^{s-1}(x^a)h_i^{s-1}(x^b)\bigr)_{a,b}
 =K(\nu_{s-1,N}).
\]
The conditional distribution depends on the previous layer only through its
empirical measure. Finally, readout weights have variance $N^{-2}$, so
\begin{equation}\label{eq:critical-output-representation}
 Z_N\ \stackrel{\mathrm{law}}{=}\ K(\nu_{r,N})^{1/2}V_N.
\end{equation}
This representation is exact for every $N$.

Sanov's theorem on $\mathcal P(\mathbb R^m)$ gives the good rate
$D(\eta\Vert\Gamma_m)$ for each $\eta_{s,N}$. The elementary Gaussian
large-deviation principle gives rate $\|v\|_2^2/2$ for $V_N$.
Independence and the finite-product principle therefore give the good rate
\[
 \sum_{s=1}^r D(\eta_s\Vert\Gamma_m)+\frac12\|v\|_2^2
\]
for $(\eta_{1,N},\ldots,\eta_{r,N},V_N)$.

The recursive map from these variables to $(\nu_1,\ldots,\nu_r,z)$ is
continuous. To verify the only potentially delicate point, suppose
$Q_j\to Q$ and $\eta_j\Rightarrow\eta$ weakly. For every continuous
$b:E\to\mathbb R$, the functions $b\circ\mathsf T_{Q_j}$ converge
uniformly on each compact subset of $\mathbb R^m$ to
$b\circ\mathsf T_Q$, since matrix square roots are continuous and $\phi$
is continuous. Their absolute values are uniformly bounded. Tightness of a
weakly convergent sequence of probability measures then gives
\[
 \int b\circ\mathsf T_{Q_j}\,d\eta_j
 \longrightarrow\int b\circ\mathsf T_Q\,d\eta.
\]
Also $\nu\mapsto K(\nu)$ is continuous because every coordinate
$h_ah_b$ is continuous and bounded on $E$. This proves continuity of the
whole recursive map, including \eqref{eq:critical-output-representation}.

Continuous contraction now gives the rate by minimizing over its preimages.
For a measurable map $T$ and a reference probability $\rho$, the entropy
contraction identity is
\begin{equation}\label{eq:entropy-contraction-identity}
 \inf_{\eta:T_\#\eta=\nu}D(\eta\Vert\rho)
 =D(\nu\Vert T_\#\rho).
\end{equation}
The lower bound is data processing. When the right side is finite and
$q=d\nu/d(T_\#\rho)$, the measure $d\eta=q(T(u))\,d\rho(u)$ attains
equality. If the right side is infinite, data processing already proves
the identity in the extended real sense. Apply
\eqref{eq:entropy-contraction-identity} successively with
$T=\mathsf T_{Q_0}$ and $T=\mathsf T_{K(\nu_{s-1})}$.
The last minimization is
\[
 \inf_{v:Q^{1/2}v=z}\frac12\|v\|_2^2=\mathcal G_Q(z),
\]
including singular $Q$. This proves \eqref{eq:critical-rate}.

The contraction theorem supplies a good rate and exponential tightness.
For completeness, the displayed minimization is also directly well behaved:
relative entropy is jointly lower semicontinuous, the map
$\nu\mapsto\mathsf M_{K(\nu)}$ is weakly continuous, and
\[
 \mathcal G_Q(z)=\sup_{t\in\mathbb R^m}
 \left\{t^{\mathsf T}z-\frac12t^{\mathsf T}Qt\right\}
\]
is jointly lower semicontinuous in $(Q,z)$. Since $\mathcal P(E)^r$ is
compact, every finite infimum in \eqref{eq:critical-rate} is attained.
Moreover $K(\nu)\preceq mI_m$, so $I(z)\geq\|z\|_2^2/(2m)$.
\end{proof}

\begin{corollary}[Full-margin masses and posterior rates]
\label{cor:critical-margin-rates}
For a Boolean function $f$, regarded as a vector in $\{-1,1\}^m$, set
\[
 J(f)=\inf_{z:\ f_a z_a\geq1/2\ \forall a}I(z).
\]
Then the strict-margin event used throughout this paper has the exact rate
\begin{equation}\label{eq:critical-full-margin-rate}
 \lim_{N\to\infty}-\frac1N\log\mu_{N,1}(A_{N,f})=J(f).
\end{equation}
For a fixed noiseless training set $S$, let $\mathcal B_S$ be the Boolean
functions agreeing with its labels. If
$J_S=\min_{g\in\mathcal B_S}J(g)<\infty$, the actual globally conditioned
posterior satisfies, for every $f\in\mathcal B_S$,
\[
 \lim_{N\to\infty}-\frac1N\log\Pi_{N,1}(f\mid S)=J(f)-J_S.
\]
In particular, posterior mass concentrates on the minimizers of $J$ within
$\mathcal B_S$. These logarithmic asymptotics do not determine how mass is
divided between tied minimizers.
\end{corollary}

\begin{proof}
For $c\geq1$, retaining any feature laws in \eqref{eq:critical-rate} and
scaling only the output gives
\[
 I(cz)\leq c^2 I(z),
\]
because the entropy terms are nonnegative and the Gaussian energy scales
quadratically. If $f_a z_a\geq1/2$ for all $a$, then $cz$ has strict
margin for every $c>1$. Letting $c\downarrow1$ shows that the infimum of
$I$ over the open strict-margin orthant equals its infimum over the
closed orthant. The lower and upper large-deviation bounds therefore agree,
proving \eqref{eq:critical-full-margin-rate}.

Conditioning the prior on the global margin event and subsequently on the
labels gives exactly
\[
 \Pi_{N,1}(f\mid S)=
 \frac{\mu_{N,1}(A_{N,f})}
 {\sum_{g\in\mathcal B_S}\mu_{N,1}(A_{N,g})}.
\]
For fixed $d$ this is a finite sum, whose exponential rate is the smallest
$J(g)$ among its terms. This proves the ratio formula and concentration
claim.
\end{proof}

The variational principle exhibits the distinction between the critical
regime and a frozen kernel model. At $\gamma=1$, deviations of a hidden
feature law have an order-$N$ entropy cost, at the same scale as an order-one
output. Thus the optimization is over every hidden feature law, with a
separate entropy charge for each layer. It is not in general legitimate to
set those laws to their unconditioned values. The following sections obtain a
polynomial description of the minimizing feature laws for $\tanh$ and use it
to compress every Boolean function with bounded full-margin rate.

\section{Quadratic tilts at a critical full-margin minimizer}
\label{sec:critical-stationarity}

We now specialize to $\phi=\tanh$. Its analyticity and its bounded first two
derivatives yield additional structure at the minimizers of
\eqref{eq:critical-rate}. The following argument uses the exact variational
principle above; it does not assume mean-field equations.

\begin{lemma}[Positive definite feature covariances]
\label{lem:critical-positive-definite}
If $(\nu_1,\ldots,\nu_r,z)$ has finite cost in
\eqref{eq:critical-rate}, then every $K(\nu_s)$ is positive definite.
\end{lemma}
\begin{proof}
For any nonzero $c\in\mathbb R^m$, the function
\[
 H_c(b,v)=\sum_{a=1}^m c_a\tanh(b+v\cdot x^a)
\]
is real analytic and is not identically zero. To see the latter, choose
$v_0$ such that the numbers $v_0\cdot x^a$ are all distinct. If $H_c$ were
identically zero, evaluating it at $(b,v)=(-t\theta,tv_0)$ and letting
$t\to\infty$ would give
\[
 \sum_a c_a\operatorname{sign}(v_0\cdot x^a-\theta)=0
\]
for every $\theta$ different from these numbers. The difference between
two successive such identities is twice one coefficient $c_a$, so all
coefficients would vanish. A nonzero real-analytic function has a
Lebesgue-null zero set. The first-layer weights have a positive Gaussian
density on the $(d+1)$-dimensional $(b,v)$ space. Consequently
\[
 \mathsf M_{Q_0}\{h:c^{\mathsf T}h=0\}=0.
\]
Finite entropy implies $\nu_1\ll\mathsf M_{Q_0}$, and therefore
$c^{\mathsf T}K(\nu_1)c>0$ for every $c\ne0$.

Inductively, if $Q=K(\nu_{s-1})$ is positive definite, then
$\mathsf M_Q$ has a positive density on $(-1,1)^m$, since coordinatewise
$\tanh$ is a smooth bijection from $\mathbb R^m$ onto that set.
Every proper linear hyperplane is $\mathsf M_Q$-null. Absolute continuity
$\nu_s\ll\mathsf M_Q$ again gives positive definiteness of $K(\nu_s)$.
\end{proof}

\begin{theorem}[Quadratic form of the optimal feature laws]
\label{thm:critical-quadratic-tilts}\label{thm:critical-tilts}
Let $a>0$, $f\in\{-1,1\}^m$, and suppose
\[
 J_a(f)=\inf_{z:\ f_bz_b\geq a\ \forall b}I(z)\leq K<\infty.
\]
There is a minimizing tuple $(\nu_1,\ldots,\nu_r,z)$.
Write $Q_s=K(\nu_s)$ and $\beta=Q_r^{-1}z$.
There are symmetric matrices $\Gamma_1,\ldots,\Gamma_r$ such that
\begin{equation}\label{eq:critical-quadratic-tilts}
 \frac{d\nu_1}{d\mathsf M_{Q_0}}(h)
   =\frac{e^{h^{\mathsf T}\Gamma_1h}}{Z_1},\qquad
 \frac{d\nu_s}{d\mathsf M_{Q_{s-1}}}(h)
   =\frac{e^{h^{\mathsf T}\Gamma_sh}}{Z_s}\quad(2\leq s\leq r).
\end{equation}
Moreover,
\begin{equation}\label{eq:critical-beta-bound}
 \|\beta\|_1\leq\frac{2K}{a},\qquad
 \Gamma_r=\frac12\beta\beta^{\mathsf T}.
\end{equation}
If $g_s=\sum_{i,j}|(\Gamma_s)_{ij}|$, then
\begin{equation}\label{eq:critical-gamma-recursion}
 g_r\leq\frac{2K^2}{a^2},\qquad
 g_s\leq3g_{s+1}+2g_{s+1}^2\quad(1\leq s<r).
\end{equation}
Thus, at fixed depth, all these entrywise norms are bounded by a polynomial
in $K$ and $1/a$, independently of the number $m=2^d$ of inputs.
\end{theorem}

\begin{proof}
Existence follows from the goodness of $I$, closedness of the margin
constraints, and attainment of the feature-law infimum in
Theorem~\ref{thm:critical-ldp}. Every $Q_s$ is positive definite by
Lemma~\ref{lem:critical-positive-definite}.

For fixed feature laws, the output minimizes the strictly convex quadratic
$z^{\mathsf T}Q_r^{-1}z/2$ subject to the linear margin constraints.
Their feasible set has a strictly feasible point. The elementary
quadratic-program optimality conditions give multipliers $\lambda_b\geq0$
such that
\[
 \beta_b=f_b\lambda_b,\qquad
 \lambda_b(f_bz_b-a)=0.
\]
It follows that
\[
 \frac12\beta^{\mathsf T}z
 =\frac a2\sum_b\lambda_b
 =\frac a2\|\beta\|_1\leq J_a(f)\leq K,
\]
which proves the first bound in \eqref{eq:critical-beta-bound}.

We justify the feature-law variations, including possible zeros of their
densities. Fix a layer $s$, freeze every other variable, and denote its
reference measure by $\mu_s$ (either $\mathsf M_{Q_0}$ or
$\mathsf M_{Q_{s-1}}$). Its contribution to the objective has the form
\[
 D(\nu_s\Vert\mu_s)+F_s(K(\nu_s))+\text{constant},
\]
where $F_s$ is differentiable near the current positive definite $Q_s$.
For $s=r$, $F_r(Q)=z^{\mathsf T}Q^{-1}z/2$, so
\[
 \nabla F_r(Q_r)=-\frac12\beta\beta^{\mathsf T}.
\]
For $s<r$, put $\xi(h)=(\operatorname{arctanh}(h_b))_{b=1}^m$ and
$B_{s+1}=\int\xi(h)\xi(h)^{\mathsf T}\,d\nu_{s+1}(h)$.
This moment is finite: by invariance of relative entropy under the
coordinatewise $\tanh$ bijection, the preactivation law has finite
relative entropy with respect to $\mathcal N(0,Q_s)$, and the Gaussian
exponential moment of a sufficiently small multiple of
$\xi^{\mathsf T}Q_s^{-1}\xi$ controls its expectation by the entropy
inequality. The Gaussian density ratio, whose Jacobian cancels under
$\tanh$, gives, near $Q_s$,
\[
 F_s(Q)=\frac12\log\det Q+\frac12\operatorname{tr}(Q^{-1}B_{s+1})
       +\text{constant}.
\]
In particular,
\begin{equation}\label{eq:critical-Gamma-from-moments}
 -\nabla F_s(Q_s)
 =\frac12Q_s^{-1}(B_{s+1}-Q_s)Q_s^{-1}=:\Gamma_s.
\end{equation}

Put $B=\nabla F_s(Q_s)$ and let $\rho_s$ have density
$e^{-h^{\mathsf T}Bh}/Z$ with respect to $\mu_s$. Its logarithm is bounded,
so $D(\nu_s\Vert\rho_s)$ is finite. For
$\nu_{s,t}=(1-t)\nu_s+t\rho_s$, convexity of relative entropy gives
\[
 D(\nu_{s,t}\Vert\rho_s)-D(\nu_s\Vert\rho_s)
 \leq-tD(\nu_s\Vert\rho_s).
\]
The identity
$D(\nu\Vert\rho_s)=D(\nu\Vert\mu_s)+\int h^{\mathsf T}Bh\,d\nu+\log Z$
and differentiability of $F_s$ imply that the full objective changes by
at most $-tD(\nu_s\Vert\rho_s)+o(t)$. Minimality forces
$D(\nu_s\Vert\rho_s)=0$. Thus $\nu_s=\rho_s$, proving
\eqref{eq:critical-quadratic-tilts} with
$\Gamma_s=-\nabla F_s(Q_s)$. This also proves
$\Gamma_r=\beta\beta^{\mathsf T}/2$.

It remains to bound the apparently ill-conditioned expression
\eqref{eq:critical-Gamma-from-moments} without using $\|Q_s^{-1}\|$.
Under $\nu_{s+1}$ the preactivation vector $\xi$ has Gaussian reference
$\mathcal N(0,Q_s)$ tilted by $e^{V(\xi)}$, where
\[
 V(\xi)=\tanh(\xi)^{\mathsf T}\Gamma_{s+1}\tanh(\xi).
\]
Twice applying Gaussian integration by parts gives
\[
 B_{s+1}-Q_s
 =Q_s\,\mathbb E_{\mathrm{tilt}}
       [\nabla^2V+(\nabla V)(\nabla V)^{\mathsf T}]\,Q_s.
\]
This identity is justified directly because $V$ and its first two
derivatives are bounded. Substitution into
\eqref{eq:critical-Gamma-from-moments} gives
\[
 \Gamma_s=\frac12\mathbb E_{\mathrm{tilt}}
       [\nabla^2V+(\nabla V)(\nabla V)^{\mathsf T}].
\]
The bounds $|\tanh|\leq1$, $|\tanh'|\leq1$, and
$|\tanh''|\leq2$ imply
\[
 \|\nabla V\|_1\leq2g_{s+1},\qquad
 \sum_{i,j}|\partial_i\partial_jV|\leq6g_{s+1}.
\]
Therefore $g_s\leq3g_{s+1}+2g_{s+1}^2$.
The final-layer bound is
$g_r=\|\beta\|_1^2/2\leq2K^2/a^2$.
Iterating the displayed quadratic recursion a fixed number of times proves
the asserted polynomial dependence.
\end{proof}

\section{Compression of bounded quadratic-tilt towers}\label{sec:critical-compression}
This section uses $\phi=\tanh$. Write $\|A\|_{\max}=\max_{ij}|A_{ij}|$ and
$\|A\|_{1,\mathrm{ent}}=\sum_{ij}|A_{ij}|$ for matrices on the finite input
set. These are different norms; their dual pairing is what avoids a factor
of $2^d$ below.
For a positive semidefinite covariance $Q$ and a symmetric matrix $\Gamma$,
let
\[
 \mathsf F_\Gamma(Q)=
 \frac{\mathbb E_{Z\sim\mathcal N(0,Q)}
 [\phi(Z)\phi(Z)^\top\exp\{\phi(Z)^\top\Gamma\phi(Z)\}]}
 {\mathbb E_{Z\sim\mathcal N(0,Q)}
 [\exp\{\phi(Z)^\top\Gamma\phi(Z)\}]}.
\]
All activations are coordinatewise. The definition is valid for singular $Q$.

\begin{lemma}[A dimension-independent covariance comparison]\label{lem:tilted-covariance}
If $g=\|\Gamma\|_{1,\mathrm{ent}}$, then for all positive semidefinite
$Q,Q'$ of the same size,
\[
 \|\mathsf F_\Gamma(Q)-\mathsf F_\Gamma(Q')\|_{\max}
 \leq 8(1+g)^2\|Q-Q'\|_{\max}.
\]
\end{lemma}
\begin{proof}
Put $V(z)=\phi(z)^\top\Gamma\phi(z)$ and
$Q_t=(1-t)Q+tQ'$. For a twice differentiable bounded test function $F$, let
$\mathbb E_t$ denote expectation under the Gaussian of covariance $Q_t$
reweighted by $e^V$. Gaussian differentiation and the quotient rule give
\begin{align*}
 \frac{d}{dt}\mathbb E_tF
 =\frac12\sum_{a,b}(Q'-Q)_{ab}\bigl\{
 &\mathbb E_t[\partial_{ab}F+
       \partial_aF\partial_bV+\partial_bF\partial_aV]\\
 &+\operatorname{Cov}_t(F,\partial_{ab}V+
                                   \partial_aV\partial_bV)\bigr\}.
\end{align*}
This formula has no reciprocal partition-function bound: all terms are
expectations under a probability measure. Define
$M_1(H)=\sup_z\sum_a|\partial_aH(z)|$ and
$M_2(H)=\sup_z\sum_{a,b}|\partial_{ab}H(z)|$.
Because $|\phi|,|\phi'|\leq1$ and $|\phi''|\leq2$,
\[
 M_1(V)\leq2g,\qquad M_2(V)\leq6g.
\]
For $F(z)=\phi(z_i)\phi(z_j)$, including $i=j$,
$\|F\|_\infty\leq1$, $M_1(F)\leq2$, and $M_2(F)\leq6$.
The absolute derivative is therefore at most
\[
 \frac12\|Q-Q'\|_{\max}
 \{6+8g+2(6g+4g^2)\}
 =(3+10g+4g^2)\|Q-Q'\|_{\max}.
\]
Integrating in $t$ proves the assertion, with the displayed slightly larger
constant. One may initially replace both covariances by $Q+\epsilon I$ and
$Q'+\epsilon I$. The functions and all derivatives used above are bounded,
and the partition functions lie between $e^{-g}$ and $e^g$, so weak Gaussian
convergence lets $\epsilon\downarrow0$. This also proves the singular case.
\end{proof}

\begin{proposition}[Finite realization of a bounded-tilt tower]\label{prop:tilt-compress}
Fix $L\geq2$, $K\geq1$, and $G\geq1$. Let $Q_0$ be the covariance of the
Gaussian first-layer field $w\cdot\bar x$, with
$w\sim\mathcal N(0,I_{d+1}/d)$, indexed by $x\in\mathcal X$.
Suppose symmetric matrices $\Gamma_1,\ldots,\Gamma_{L-1}$ and a vector
$\beta\in\mathbb R^{\mathcal X}$ satisfy
\[
 \|\Gamma_s\|_{1,\mathrm{ent}}\leq G,\qquad
 \|\beta\|_1\leq4K,
 \qquad Q_s=\mathsf F_{\Gamma_s}(Q_{s-1}).
\]
Then the function $z(x)=(Q_{L-1}\beta)_x$ can be uniformly approximated to
accuracy $1/8$ by a depth-$L$ network whose width and weights are bounded by
$\operatorname{poly}_L(d,K,G)$.
If $f(x)z(x)\geq1/2$ everywhere, a factor-two increase in the finite
network's readout implements $f$ with margin at least $3/4$.
\end{proposition}
\begin{proof}
Let $b=8(1+G)^2$ and
\[
 \eta=[32K(L-1)b^{L-2}]^{-1}.
\]
Choose the common width $n$ to be the least integer satisfying
\[
 n\geq2\eta^{-2}\log\bigl(8(L-1)\,2^{2d}\bigr).
\]
The construction is random only to prove existence.
The matrices $\Gamma_s$ may be indexed by the entire cube; they select the
sampling laws, but are not stored or evaluated by the resulting finite
network. No efficient sampling or construction algorithm is claimed.
At the first layer, independently sample $n$ Gaussian weight vectors
$w_i\sim\mathcal N(0,I_{d+1}/d)$ tilted by
\[
 \exp\{h(w_i)^\top\Gamma_1h(w_i)\},\qquad
 h(w_i)_x=\phi(w_i\cdot\bar x).
\]
Their activation vectors have exactly the law defining
$Q_1=\mathsf F_{\Gamma_1}(Q_0)$.
At each subsequent hidden layer, conditional on the already sampled previous
activation vectors $h_1,\ldots,h_n$, independently sample $n$ standard
Gaussian row vectors $\xi_j\in\mathbb R^n$ tilted by
\[
 \exp\{\phi(Z_j)^\top\Gamma_s\phi(Z_j)\},
 \qquad Z_j(x)=\frac1{\sqrt n}\sum_{i=1}^n\xi_{ji}h_i(x).
\]
Use the actual network weights $\xi_{ji}/\sqrt n$.
If $\widehat Q_{s-1}=n^{-1}\sum_i h_ih_i^\top$ is the preceding empirical
activation covariance, then, conditionally, $Z_j$ before tilting is Gaussian
with exactly that covariance. Thus the new activation vectors are
conditionally independent, with second moment
$\mathsf F_{\Gamma_s}(\widehat Q_{s-1})$.

Every covariance entry is the empirical mean of $n$ variables in $[-1,1]$.
Hoeffding's inequality, a union bound over $2^{2d}$ entries, and conditioning
at successive layers show, with probability at least $3/4$, that
\[
 \|\widehat Q_s-
          \mathsf F_{\Gamma_s}(\widehat Q_{s-1})\|_{\max}\leq\eta
 \quad(1\leq s<L),
\]
where $\widehat Q_0=Q_0$.
Lemma~\ref{lem:tilted-covariance} then implies
\[
 \|\widehat Q_{L-1}-Q_{L-1}\|_{\max}
 \leq (L-1)b^{L-2}\eta.
\]
For each last-layer activation vector $h_i$, choose the deterministic
readout weight $v_i=(\beta\cdot h_i)/n$. The resulting output is
$\widehat z(x)=(\widehat Q_{L-1}\beta)_x$, and therefore
\[
 \|\widehat z-z\|_\infty
 \leq4K(L-1)b^{L-2}\eta=1/8.
\]
Although $\beta$ is indexed by the entire input cube, its only role in the
finite network is to select these $n$ real readout weights; each satisfies
$|v_i|\leq4K/n$. It is not an additional computational input to the network.

It remains to bound the sampled weights without losing polynomial size.
For every sampled row, its tilted density relative to its Gaussian reference
is at most $e^{2G}$, since the quadratic tilt lies in $[-G,G]$.
Let $D_n=n(d+1)+(L-2)n^2$ be the total number of incoming scalar coordinates,
and choose
\[
 B=\bigl[4G+2\log(8\max\{D_n,1\})\bigr]^{1/2}.
\]
Each raw incoming coordinate has reference variance at most one, so, also
conditionally on all preceding layers, its probability of exceeding $B$ in
absolute value is at most $2e^{2G-B^2/2}$.
A union bound gives probability at most $1/4$ that any incoming coordinate
exceeds $B$. The actual internal weights divide the raw coordinates by
$\sqrt n$ and so satisfy the same bound. Intersecting this event with the
covariance-approximation event leaves probability at least $1/2$.
Fix a realization in the intersection. Both $n$ and
$\max\{B,8K/n\}$ after the factor-two readout increase are bounded by
polynomials in $d,K,G$ at fixed depth. The stated margin follows from
$2(1/2-1/8)=3/4$.
\end{proof}

\begin{lemma}[Counting finite realizations]\label{lem:tilt-count}
Fix $L$, width $n$, and an entrywise bound $B\geq1$. The Boolean functions
having depth-$L$, width-$n$ implementations with margin at least $3/4$ and
weights bounded by $B$ number at most
\[
 \exp\{P\log(3+16B\Lambda)\},\qquad
 P=n(d+1)+(L-2)n^2+n,
 \quad \Lambda=(d+Ln)[1+n(B+1)]^{L-1}.
\]
\end{lemma}
\begin{proof}
Round every weight to a multiple of $(8\Lambda)^{-1}$.
The first hidden layer's coordinatewise error is at most $(d+1)/(8\Lambda)$.
Every subsequent hidden layer has error at most $n(B+1)$ times the preceding
error plus $n/(8\Lambda)$; the same estimate applies at the linear readout.
Induction bounds the final error by $1/8$. Each rounded coordinate has at
most $3+16B\Lambda$ possible values, so the $P$-parameter grid has the stated
size. Rounding preserves every represented Boolean sign function.
\end{proof}

\section{Critical rate complexity and Bayesian generalization}
\begin{theorem}[Polynomial rate-sublevel counting]\label{thm:critical-counting}
For every fixed depth $L\geq2$, with activation $\tanh$, there is a polynomial
$B_L$ such that, for every $K\geq1$,
\[
 \log\#\{f:\mathcal X\to\{-1,1\}:J(f)\leq K\}\leq B_L(d,K).
\]
Moreover every function in this sublevel set has a depth-$L$ implementation
of width and weight magnitude at most $\operatorname{poly}_L(d,K)$.
\end{theorem}
\begin{proof}
Take a minimizing tuple for $J(f)\leq K$ in
Theorem~\ref{thm:critical-tilts}, with margin $a=1/2$.
Its vector $\beta$ satisfies $\|\beta\|_1\leq4K$.
Let $g_r^*=8K^2$ and recursively
$g_s^*=3g_{s+1}^*+2(g_{s+1}^*)^2$ for $s<r=L-1$.
Then $G=1+\max_s g_s^*$ is a polynomial in $K$ at fixed depth,
and all the quadratic tilt matrices have entrywise norm at most $G$.
The feature covariances satisfy
$Q_s=\mathsf F_{\Gamma_s}(Q_{s-1})$, and the optimal output is
$z=Q_r\beta$.
Proposition~\ref{prop:tilt-compress} supplies a finite implementation with
margin at least $3/4$, common width bounded by a polynomial in $d,K$, and
weights bounded by another such polynomial.
The bounds in its construction can be chosen uniformly for every function
in the sublevel set. Lemma~\ref{lem:tilt-count} counts the resulting fixed
parameter grid and gives the stated bound.
For the complexity statement, increase width by zero padding if necessary
so that the width also bounds every weight magnitude. This preserves the
network for $\tanh$, since $\tanh(0)=0$.
\end{proof}

\begin{proof}[Proof of Theorem~\ref{thm:critical-learning}]
The uniform cloning bound and the exact margin rate imply
\[
 J(f)\leq Q_L(d,C)/2,\qquad C=C_L(f).
\]
Let
\[
 \mathcal A_f=\{g:J(g)\leq J(f)+1\}.
\]
Its logarithmic cardinality is bounded by
$B_L(d,Q_L(d,C)/2+1)$, and hence by a polynomial $p_L(d,C)$.
For each $g\notin\mathcal A_f$, the exact rate formula gives
\[
 \frac{\mu_{N,1}(A_{N,g})}{\mu_{N,1}(A_{N,f})}\longrightarrow0
 \quad(N\to\infty).
\]
Indeed the rate gap is greater than one; an infinite rate causes no
problem. At fixed $d$ there are finitely many Boolean functions, so
summing these ratios yields
\[
 \frac{\pi_{N,1}(\mathcal A_f^c)}{\pi_{N,1}(f)}\longrightarrow0.
\]
Thus there is $N_0\geq C$ such that for every $N\geq N_0$ this ratio is
at most $\varepsilon/6$. This width threshold depends on the target and
accuracy, but not on the input distribution or training sample.
Lemma~\ref{lem:relevant}, with $a=\tau=\varepsilon/6$ and
$H=p_L(d,C)$, now gives $R_D(\Pi_{N,1})\leq\varepsilon/3$ with
probability at least $1-\delta$ whenever
\[
 n\geq\frac6\varepsilon[p_L(d,C)+\log(1/\delta)].
\]
At every input where the correct-label posterior probability is at most
$2/3$, the posterior error is at least $1/3$. The distributional mass of
such inputs is at most $3R_D(\Pi_{N,1})\leq\varepsilon$.
This proves the claimed confidence guarantee, including its strict
success threshold.
\end{proof}


\section{Shallow norm compression and counting}
\begin{lemma}[Explicit shallow counting]\label{lem:count}
For $T\geq1$, let $\mathcal A_{N,T}$ be the set of Boolean functions
implemented with strict signed margin greater than $1/2$ at every input
and $\kappa_N(w)\leq T$ by a width-$N$ network with $L=2$.
Let
\[
 m=\left\lceil512(T+1)((d+1)\log2+1)\right\rceil,
 \quad
 H(d,T)=m(d+2)\log\bigl(2^{14}m(d+1)(T+1)^2\bigr).
\]
Then $\log|\mathcal A_{N,T}|\leq H(d,T)$, independently of $N$.
\end{lemma}
\begin{proof}
Write the incoming rows as $u_i\in\R^{d+1}$ and the readout entries as $v_i$.
The norm bound implies
\[
 \sum_i\|u_i\|_2^2\leq NT/d,\quad
 \sum_i v_i^2\leq T/N,\quad
 \sum_i|v_i|\leq\sqrt T.
\]
Delete terms with $\|u_i\|_2>B$, where $B=16T/\sqrt d$.
There are at most $NT/(dB^2)$ such terms. Cauchy--Schwarz bounds their
uniform contribution to the output by $T/(\sqrt d B)=1/16$.
For the retained terms put $S=\sum_i|v_i|\leq\sqrt T$.
An implementing network cannot have $S=0$ after this deletion, since its
original absolute output was everywhere greater than $1/2$.
Draw $m$ retained indices independently with probabilities $|v_i|/S$.
The empirical network
\[
 \widetilde z(x)=\frac S m\sum_{j=1}^m
      \operatorname{sgn}(v_{i_j})\phi(u_{i_j}\cdot\bar x)
\]
has expectation equal to the retained network. For each input, Hoeffding's
inequality gives
\[
 \Prb\{|\widetilde z(x)-\E\widetilde z(x)|>1/16\}
 \leq2\exp[-m/(512T)].
\]
A union bound over the $2^d$ inputs makes the failure probability less than one
for the stated $m$. Fix a realization with uniform error at most $1/16$.
The total error from the original network is at most $1/8$.
Doubling the readout now gives a width-$m$ implementation with margin
strictly greater than $3/4$, incoming coordinates bounded by $16T$, and
readout absolute sum at most $2\sqrt T$.

Round each incoming coordinate to the nearest multiple of
$\eta_1=[64\sqrt{T+1}(d+1)]^{-1}$ and each readout coordinate to the nearest
multiple of $\eta_2=(32m)^{-1}$. Rounding incoming weights changes the output
by at most $2\sqrt T(d+1)\eta_1\leq1/32$; rounding the readout changes it by
at most $m\eta_2=1/32$. These conservative bounds allow rounding errors of a
whole mesh step. The resulting grid network preserves the sign of the
implementation at every input.
Each incoming coordinate has at most $3+2048(d+1)(T+1)^{3/2}$ grid choices.
Each readout coordinate, bounded before rounding by $2\sqrt T$, has at most
$3+128m\sqrt{T+1}$ grid choices. Their product, over $m(d+1)$ incoming and
$m$ readout coordinates, is at most $\exp H(d,T)$.
Every function in the class has a sign-preserving representative on this one
fixed grid. A grid network computes at most one Boolean sign function, so the
same bound counts the functions themselves.
\end{proof}

\section{The posterior bound}
\begin{lemma}[Gaussian tail]
If $X\sim\chi^2_D$ and $u\geq8D$, then $\Prb(X\geq u)\leq e^{-u/4}$.
\end{lemma}
\begin{proof}
Chernoff's inequality at parameter $3/8$ gives
$\Prb(X\geq u)\leq4^{D/2}e^{-3u/8}$.
Since $D\leq u/8$, its logarithm is at most
$u\log2/8-3u/8\leq-u/4$.
\end{proof}
\begin{lemma}[Finite relevant class]\label{lem:relevant}
Let a prior on a finite class assign positive mass to a target $f$. Suppose
$f\in\mathcal A$, $|\mathcal A|\leq e^H$, and
$\pi(\mathcal A^c)\leq\tau\pi(f)$.
For any $a,\delta\in(0,1)$, if
$n\geq a^{-1}(H+\log(1/\delta))$, then with probability at least
$1-\delta$ its consistency posterior satisfies $R_D(\Pi)\leq a+\tau$.
\end{lemma}
\begin{proof}
Any fixed function with error greater than $a$ survives $n$ labels with
probability at most $(1-a)^n\leq e^{-an}$. A union bound excludes all such
survivors in $\mathcal A$ except on an event of probability $e^{H-an}$.
For every sample, the posterior normalizing probability is at least
$\pi(f)$, so its mass outside $\mathcal A$ is at most $\tau$.
\end{proof}
\begin{proof}[Proof of Theorem~\ref{thm:shallow}]
Let $Q=Q_2(d,C)$ be the explicit cloning bound above, so
$\mu_{N,\gamma}(A_{N,f})\geq e^{-N^\gamma Q/2}$.
Choose
\[
 T=\max\{1,\ 8(d+2),\ 2Q+4\log(2/\varepsilon)\}.
\]
At $L=2$ the number of Gaussian parameters is $D_0=N(d+2)$ and
$N^\gamma\kappa_N(w)\sim\chi^2_{D_0}$.
Since $\gamma\geq1$, $D_0\leq(d+2)N^\gamma$, so the preceding Gaussian
bound yields
\[
 \frac{\mu_{N,\gamma}\{\kappa_N>T\}}
      {\mu_{N,\gamma}(A_{N,f})}
 \leq\exp[-N^\gamma(T/4-Q/2)]\leq\varepsilon/2.
\]
In particular $A_{N,f}$ contains a point with norm at most $T$, and thus
$f\in\mathcal A_{N,T}$. If a function is outside $\mathcal A_{N,T}$, every
one of its implementing weights has norm greater than $T$. Disjointness of
the implementing events and cancellation of the global-margin normalization
therefore give
\[
 \frac{\pi_{N,\gamma}(\mathcal A_{N,T}^{\,c})}{\pi_{N,\gamma}(f)}
 \leq\varepsilon/2.
\]
Apply Lemmas~\ref{lem:count} and~\ref{lem:relevant} with
$a=\tau=\varepsilon/2$ and $H=H(d,T)$.
The explicit quantity $H(d,T)$ is bounded by a polynomial in
$d,C,\log(2/\varepsilon)$, proving the theorem.
For posterior majority, every misclassified input has posterior probability
of error at least $1/2$; integration gives its error at most $2R_D(\Pi)$.
\end{proof}

\section{Scope and relation to earlier mean-field arguments}
At $\gamma=1$ and depth $L\geq3$, the typical weighted squared norm is of
order $N$ because there are $(L-2)N^2$ internal weights. The shallow
norm-tail argument therefore cannot be used directly. Our all-depth proof
instead bounds the coefficients of the optimal feature-law tilts and samples
finite networks from their tilted Gaussian rows. The strong-prior parameter
$\gamma$ in the common cloning construction does not enter that endpoint
variational argument.

The all-depth theorem is asymptotic in width, at fixed input dimension and
depth. Its sample polynomial has width-independent coefficients, while its
width threshold follows from a finite union of large-deviation comparisons
and is not quantified. The shallow theorem holds at every width $N\geq C$.
For larger depth the arguments in this paper do not establish an all-width
endpoint theorem, a polynomial bound on the width threshold, or the same
all-depth conclusion for
every bounded Lipschitz activation. Tanh supplies analyticity, invertibility,
and bounded first and second derivatives at specific steps of our proof.

Vaintrob and H\"anni~\cite{mft} propose an all-depth mean-field compression
argument using a finite-dimensional active subspace and fresh Gaussian
fluctuations in its complement. There are two independent reasons to supply a separate argument
for the all-depth conclusion here. First,
its learning rule uses a smooth mean-squared-error Gibbs tilt, and its
architecture has independent biases at every hidden layer. Our prior instead
has a global Boolean margin restriction followed by exact label conditioning.
Second, the displayed regularity invariant used in its proof is false as
stated. We give a direct check, rather than assume that a change of conventions
would repair it.

\begin{proposition}[Failure of the stated pointwise sensitivity invariant]
Let $\mu$ be uniform on $\mathcal X$, $M\geq1$, and
$S(\xi)(x)=M^{-1}\sum_{j=1}^M\tanh(\xi_j(x))$.
If the bound
\[
 |D_{j,u}S(\xi)(x)|\leq\frac A{\sqrt M}\|u\|_{L^2(\mu)}
 \tag{$*$}
\]
holds for every field $\xi$, input $x$, channel $j$, and direction
$u\in L^2(\mu)$, then $A\geq2^{d/2}/\sqrt M$.
\end{proposition}
\begin{proof}
Set $\xi=0$ and $u=2^{d/2}{\bf1}_{\{x\}}$. Then
$\|u\|_{L^2(\mu)}=1$, whereas
$D_{j,u}S(0)(x)=2^{d/2}/M$. Substitution into $(*)$ gives the result.
\end{proof}
The invariant in Appendix A, Definition A.6(I4), of~\cite{mft} has precisely
the form $(*)$. Appendix B, Theorem B.14, asserts its initialization for the
readout suffix with polynomial width and polynomial regularity constant.
The proposition contradicts that assertion for polynomial $M$ as $d$ grows.
A pointwise derivative estimate involving $|u(x)|$, or an averaged estimate
for population loss, can be true. Replacing $(*)$ by either one changes what
has to be proved in the subsequent induction, particularly for the empirical
loss and conditional lazy-coordinate scores.

This identifies a gap in that earlier argument, not a counterexample to
mean-field learnability. The present proof does not use the false sensitivity
invariant, a conditional lazy-score bound, or a propagation-of-chaos premise.
Its exact finite-width distributional representation and Gaussian covariance
comparison yield the required count directly.

\begin{thebibliography}{9}
\bibitem{dembozeitouni} Amir Dembo and Ofer Zeitouni.
\emph{Large Deviations Techniques and Applications}. Second edition.
Springer, 1998. Sanov's theorem and the contraction principle.
\bibitem{mft} Dmitry Vaintrob and Kaarel H\"anni.
\emph{Width-Robust Learnability in Mean-Field Bayesian Neural Networks}.
2026. \href{https://arxiv.org/abs/2607.05735}{arXiv:2607.05735v1}.
\bibitem{slides} Dmitry Vaintrob and Kaarel H\"anni.
\emph{A Generalization Result for Overparametrized Neural Net Bayesian Learning}.
Lecture notes, 2026.
\url{https://kaarelh.github.io/ai/nnbayes-generalization.pdf}.
\end{thebibliography}
\end{document}
